% Commandes personnalisées réutilisables pour le CV

% === COMMANDE EXPÉRIENCE ===
% Usage: \experience{entreprise}{poste}{lieu}{dates}{description}{réalisations}
\newcommand{\experience}[6]{
    \subsection{#2}
        \begin{tabularx}{\linewidth}{@{}X X@{}}
        \textbf{#1} & \hfill \textcolor{lightgray}{#3} \\
        \textcolor{emerald}{\textbf{#4}} & \\
    \end{tabularx}
    \vspace{2pt}
    \textit{#5}
    \vspace{2pt}
    \begin{itemize}
        #6
    \end{itemize}
    \vspace{6pt}
}

% === COMMANDE PROJET ===
% Usage: \project{nom}{dates}{site}{technologies}{description}{réalisations}
\newcommand{\project}[6]{
    \subsection{#1}
    \begin{tabularx}{\linewidth}{@{}X@{}}
        \ifXeTeX
            \textcolor{emerald}{\faCalendar\ #2} & \hfill
            \href{#3}{\faGlobe\ Site}
        \else
            \textcolor{emerald}{\ding{182}\ #2} & \hfill
            \href{#3}{Site}
        \fi
    \end{tabularx}
    \vspace{2pt}
    \textbf{Stack technique :} \tech{#4}
    \vspace{2pt}
    \textit{#5}
    \vspace{2pt}
    \begin{itemize}
        #6
    \end{itemize}
    \vspace{8pt}
}

% === COMMANDE COMPÉTENCE ===
% Usage: \skill{catégorie}{liste des compétences}
\newcommand{\skill}[2]{
    \textbf{#1 :} \tech{#2}
    \vspace{3pt}
}

% === COMMANDE FORMATION ===
% Usage: \education{diplôme}{établissement}{dates}{mention}
\newcommand{\education}[4]{
    \textbf{#1} \\
    \textcolor{emerald}{#2} \hfill \textcolor{lightgray}{#3} \\
    \ifx&#4&
    \else
        \textit{#4}
    \fi
    \vspace{6pt}
}

% === COMMANDE CERTIFICATION ===
% Usage: \certification{nom}{organisme}{date}
\newcommand{\certification}[3]{
    \textbf{#1} \\
    \textcolor{emerald}{#2} \hfill \textcolor{lightgray}{#3}
    \vspace{3pt}
}

% === COMMANDE LANGUE ===
% Usage: \language{langue}{niveau}
\newcommand{\language}[2]{
    \textbf{#1} : \tech{#2}
    \vspace{3pt}
}

% === COMMANDE CONTACT ===
% Usage: \contact{icône}{texte}{lien}
\newcommand{\contact}[3]{
    \ifx&#3&
        \fa#1\ #2
    \else
        \href{#3}{\fa#1\ #2}
    \fi
}

% === COMMANDE LISTE RÉALISATIONS ===
% Usage: \achievement{réalisation}
\newcommand{\achievement}[1]{
    \item #1
}

% === COMMANDE INTÉRÊT ===
% Usage: \interest{intérêt}
\newcommand{\interest}[1]{
    \item #1
}
